\chapter{Compact Generation of \(\mathrm{IndCoh}_{\mathrm{Nilp}}(\mathcal{X}^{\mathrm{EG},p}_{GL_{2}})\) at Wild Primes via Residual Gerbes and Serre Twists}
\label{v4-arith:w5-g:f1a3-C}

The conjecture of Chapter~\ref{v4-arith:f1a3-BH} (item F1.a.3.C) asks whether
\(\mathrm{IndCoh}_{\mathrm{Nilp}}(\mathcal{X}^{\mathrm{EG},p,\chi}_{GL_{2}})\)
is compactly generated at the wild primes \(p\in\{2,3\}\). The answer is
not obtained here unconditionally. Emerton--Gee arXiv:1908.07185 constructs
\(\mathcal{X}^{\mathrm{EG},p}_{GL_{2}}\) as a Noetherian formal algebraic stack
at every prime. That does not give the finite Nilp stratification,
the wild-prime quiver local models, or the residual-gerbe generator
statement used below. The chapter proves the formal implication:
if these three geometric inputs are supplied, then compact generation
follows by recollement. Without them, F1.a.3.C remains open at
\(p=2,3\).

\section{Locally Noetherian Structure of \(\mathcal{X}^{\mathrm{EG},p}_{GL_{2}}\)}
\label{v4-arith:w5-g:f1a3-C:locally-noetherian}

\subsection{Setup and Conventions}

Fix a prime \(p\) and a central character
\(\chi\colon\mathbb{Z}_{p}^{\times}\to k^{\times}\) with
\(k=\overline{\mathbb{Q}}_{\ell}\), \(\ell\neq p\). Write
\(\mathcal{X}:=\mathcal{X}^{\mathrm{EG},p,\chi}_{GL_{2}}\) for the
Emerton--Gee stack of rank-\(2\) \((\varphi,\Gamma)\)-modules with fixed
determinant compatible with \(\chi\) (Emerton--Gee arXiv:1908.07185,
\S 3). We work with the classical non-derived-enhanced stack
throughout this chapter: the derived enhancement
\Cref{v4-arith:deg-f1a:conj:EG-wild-enhancement} is not invoked here.
What is needed for F1.a.3.C is compact generation of
\(\mathrm{IndCoh}_{\mathrm{Nilp}}(\mathcal{X})\) in its classical incarnation.

\begin{proposition}[Emerton--Gee: \(\mathcal{X}\) is Noetherian formal algebraic]
\label{v4-arith:w5-g:f1a3-C:prop:EG-noetherian}
\ClaimStatusProvedElsewhere. For every prime \(p\) (including
\(p\in\{2,3\}\)) and every central character \(\chi\), the
Emerton--Gee stack \(\mathcal{X}^{\mathrm{EG},p,\chi}_{GL_{2}}\) is a
Noetherian formal algebraic stack, locally of finite type over
\(\mathrm{Spf}(\mathcal{O}_{K})\) for a finite extension \(K/\mathbb{Q}_{p}\)
containing a primitive \(p\)-th root of unity. Its underlying reduced
stack \(\mathcal{X}^{\mathrm{red}}\) is a Noetherian algebraic stack of
finite Krull dimension.
\end{proposition}

\begin{proof}[Source]
Emerton--Gee arXiv:1908.07185 Theorem 3.2.1 (algebraicity),
Theorem 4.8.4 (locally of finite type), and Corollary 3.4.7
(Noetherian reduced fibre). The quoted results hold for every prime
\(p\geq 2\): the proofs do not invoke \(p\geq 5\) or any tame-prime
hypothesis. Emerton--Gee uses Breuil's \((\varphi,\Gamma)\)-module
classification and Kisin's deformation rings, both available at every
prime.
\end{proof}

\begin{remark}[formal-algebraic output of Emerton--Gee at wild primes]
\label{v4-arith:w5-g:f1a3-C:rem:EG-output-wild}
Emerton--Gee arXiv:1908.07185 establishes, uniformly in \(p\), a Noetherian
formal algebraic stack structure on \(\mathcal{X}^{\mathrm{EG},p,\chi}_{GL_{2}}\).
At \(p\in\{2,3\}\) the stack is not smooth: the extra complexity enters through
the stratification by Serre weight and the extra connected components of the
irregular locus. Smoothness and a global resolution are separate questions;
the Noetherian property is what compact generation requires.
\end{remark}

\subsection{Quasi-Excellence via Bhatt--Scholze Prisms}

\begin{proposition}[quasi-excellence of the Emerton--Gee base]
\label{v4-arith:w5-g:f1a3-C:prop:quasi-excellent}
\ClaimStatusProvedHere. The underlying formal scheme
\(\mathrm{Spf}(\mathcal{O}_{K})\) is quasi-excellent in the sense of
Grothendieck (EGA IV.7.8). For every prime \(p\) and every closed point
\(x\in\mathcal{X}^{\mathrm{EG},p,\chi}_{GL_{2}}\), the complete local
ring \(\widehat{\mathcal{O}}_{\mathcal{X},x}\) is a Noetherian excellent
ring, of finite Krull dimension.
\end{proposition}

\begin{proof}
The base ring \(\mathcal{O}_{K}\) is a complete discrete valuation ring
of mixed characteristic; complete DVRs are excellent (Matsumura 1989
\emph{Commutative Ring Theory} Theorem 32.4). Quasi-excellence is
preserved under formal completion at a closed point by Bhatt--Scholze
arXiv:1905.08229 Proposition 3.8 (which records the standard fact for
formal neighbourhoods of closed points on Noetherian excellent rings).
The complete local ring at a closed point \(x\in\mathcal{X}\) is a
power-series ring over \(\mathcal{O}_{K}\) modulo a Noetherian ideal of
height equal to the codimension of the orbit of \(x\); by Matsumura
Theorem 32.2, completion of an excellent ring at a maximal ideal is
excellent; hence \(\widehat{\mathcal{O}}_{\mathcal{X},x}\) is excellent.
Krull dimension bounded: Emerton--Gee arXiv:1908.07185 Corollary 3.4.7
bounds \(\dim\mathcal{X}^{\mathrm{red}}\leq 5\) at each closed point
for \(GL_{2}\).
\end{proof}

\begin{corollary}[locally Noetherian with finite-dimensional complete local rings]
\label{v4-arith:w5-g:f1a3-C:cor:locally-noetherian-finite-dim}
\ClaimStatusProvedHere. \(\mathcal{X}^{\mathrm{EG},p,\chi}_{GL_{2}}\) is
locally Noetherian, of bounded Krull dimension at each closed point,
quasi-excellent as a formal algebraic stack.
\end{corollary}

\begin{proof}
Combine \Cref{v4-arith:w5-g:f1a3-C:prop:EG-noetherian} (Noetherian)
with \Cref{v4-arith:w5-g:f1a3-C:prop:quasi-excellent} (excellent
complete local rings).
\end{proof}

\subsection{Nilpotent Singular Support in the Arinkin--Gaitsgory Sense}

\begin{definition}[\(\mathrm{IndCoh}_{\mathrm{Nilp}}\) on
\(\mathcal{X}^{\mathrm{EG}}\)]
\label{v4-arith:w5-g:f1a3-C:def:IndCoh-Nilp}
\ClaimStatusDefinitional. Following Arinkin--Gaitsgory arXiv:1201.6343
and Gaitsgory--Rozenblyum DAG Vol.~II Ch.~4, the
\(\infty\)-category \(\mathrm{IndCoh}_{\mathrm{Nilp}}(\mathcal{X})\)
consists of ind-coherent sheaves on \(\mathcal{X}\) whose singular
support lies in the nilpotent cone
\(\mathrm{Nilp}\subset T^{*}\mathcal{X}\), where the nilpotent cone is
cut out by the vanishing of the regular Galois-deformation cocycle
(explicit: the image of the Galois \(H^{1}\)-cocycle in the symmetric
algebra on the cotangent complex). On a Noetherian quasi-excellent
formal algebraic stack, \(\mathrm{IndCoh}_{\mathrm{Nilp}}(\mathcal{X})\)
is a well-defined presentable stable \(\infty\)-category, and embeds as
a full subcategory of \(\mathrm{IndCoh}(\mathcal{X})\).
\end{definition}

\section{Nilp Cone Decomposition at \(p=2,3\)}
\label{v4-arith:w5-g:f1a3-C:nilp-decomposition}

\subsection{The Stratification by Residual Reducibility}

\begin{proposition}[Nilp cone stratification]
\label{v4-arith:w5-g:f1a3-C:prop:nilp-stratification}
\ClaimStatusProvedHereConditional \emph{(conditional on a wild-prime
local-model theorem for the Emerton--Gee stack).} Fix \(p\in\{2,3\}\)
and a central character
\(\chi\). The nilpotent cone inside the Emerton--Gee stack admits a
finite stratification by locally closed Noetherian substacks:
\begin{equation}
\label{v4-arith:w5-g:f1a3-C:eq:stratification}
\mathcal{X}^{\mathrm{EG},p,\chi}_{GL_{2}}\cap\mathrm{Nilp}
\;=\;\mathcal{X}^{\mathrm{irr}}\;\sqcup\;\mathcal{X}^{\mathrm{red}}\;\sqcup\;\mathcal{X}^{\mathrm{wild},p},
\end{equation}
where:
\begin{enumerate}
\item \(\mathcal{X}^{\mathrm{irr}}\) is the open substack of irreducible
residual representations (smooth; tame and wild supercuspidal types
live here).
\item \(\mathcal{X}^{\mathrm{red}}\) is the closed substack of strictly
reducible non-semisimple residual representations (the ordinary
principal-series locus).
\item \(\mathcal{X}^{\mathrm{wild},p}\) is the additional wild-prime
substack, irreducible at \(p=2\) (the \emph{irregular supersingular
locus}; the component carrying the single Ext\({}^{1}\)-class of
\Cref{v4-arith:bzn-rcompletion:prop:R4-Ext1-p2}) and reducible at
\(p=3\) into three sub-strata \(\mathcal{X}^{\mathrm{wild},3}_{L_{0}}\),
\(\mathcal{X}^{\mathrm{wild},3}_{L_{1}}\), \(\mathcal{X}^{\mathrm{wild},3}_{L_{2}}\)
(indexed by the three restricted-weight simples
\(L_{0},L_{1},L_{2}\) of \Cref{v4-arith:bzn-rcompletion:prop:R4-Ext1-p3}).
\end{enumerate}
Each stratum is a locally closed Noetherian substack of
\(\mathcal{X}^{\mathrm{EG},p,\chi}_{GL_{2}}\).
\end{proposition}

\begin{proof}
Emerton--Gee gives the ambient Noetherian formal algebraic stack and
the residual-representation stratification. It does not, by itself,
identify the Nilp cone with the finite Bushnell--Henniart/Serre-weight
list displayed above. The displayed decomposition is therefore the
wild-prime local-model input. If that input holds, each listed piece is
locally closed and Noetherian, and the displayed union is a finite
stratification by assumption.
\end{proof}

\begin{remark}[relation to the Bushnell--Henniart components of
Chapter~\ref{v4-arith:f1a3-BH}]
\label{v4-arith:w5-g:f1a3-C:rem:BH-correspondence}
The stratification
(\ref{v4-arith:w5-g:f1a3-C:eq:stratification}) refines and matches
the BH-component decomposition of
Chapter~\ref{v4-arith:f1a3-BH}: the types of
\Cref{v4-arith:f1a3-BH:prop:BH-types-p2} (unramified principal
series, tame supercuspidal, wild supercuspidal) distribute across
the strata \(\mathcal{X}^{\mathrm{red}},\mathcal{X}^{\mathrm{irr}},\mathcal{X}^{\mathrm{wild},2}\)
respectively. At \(p=3\) the four BH type classes of
\Cref{v4-arith:f1a3-BH:prop:BH-types-p3} distribute among
\(\mathcal{X}^{\mathrm{red}},\mathcal{X}^{\mathrm{irr}}\), and the three
\(\mathcal{X}^{\mathrm{wild},3}_{L_{i}}\), each principal-block BH type
corresponding to exactly one \(L_{i}\)-stratum. The stratification
introduced here is the spectral-side analogue of the BH block
decomposition.
\end{remark}

\subsection{Residual Gerbes at Closed Points of Each Stratum}

\begin{proposition}[residual gerbe at a closed point]
\label{v4-arith:w5-g:f1a3-C:prop:residual-gerbe}
\ClaimStatusProvedHere. Let
\(x\in\mathcal{X}^{\mathrm{EG},p,\chi}_{GL_{2}}\) be a closed point.
The residual gerbe \(\mathcal{G}_{x}\) at \(x\) is the closed substack
defined by the maximal ideal of the complete local ring
\(\widehat{\mathcal{O}}_{\mathcal{X},x}\); it is a Noetherian algebraic
gerbe of the form \(B\mathrm{Aut}(x)\) over the residue field
\(k(x)\). For \(\mathcal{X}^{\mathrm{EG}}\) the automorphism group is
the scheme-theoretic stabiliser of the residual Galois representation,
which is an affine algebraic group of finite type over \(k(x)\).
\end{proposition}

\begin{proof}
The definition of the residual gerbe is classical
(Laumon--Moret-Bailly 2000 \emph{Champs Alg\'ebriques} \S 11,
Stacks Project Tag 06MA). For \(\mathcal{X}^{\mathrm{EG}}\) the
automorphism scheme at a residual Galois representation \(\rho\) is
the centraliser of \(\rho\) in \(GL_{2}\) over the residue field;
concretely, this is either \(\mathbb{G}_{m}\rtimes\{\pm 1\}\) or
\(\mathbb{G}_{m}\times\mathbb{G}_{m}\) depending on whether \(\rho\) is
non-split or split, both affine of finite type over \(k(x)\). Their
smoothness over \(k(x)\) follows from Emerton--Gee arXiv:1908.07185
\S 3.4 (smooth atlas structure).
\end{proof}

\section{Explicit Compact Generators per Stratum}
\label{v4-arith:w5-g:f1a3-C:compact-generators}

\subsection{The Irreducible Stratum}

\begin{proposition}[compact generator on the irreducible stratum]
\label{v4-arith:w5-g:f1a3-C:prop:compact-gen-irr}
\ClaimStatusProvedHereConditional \emph{(conditional on the stratum
being QCA and on the closed-gerbe generation theorem in this setting).}
On the stratum \(\mathcal{X}^{\mathrm{irr}}\)
the \(\infty\)-category
\(\mathrm{IndCoh}_{\mathrm{Nilp}}(\mathcal{X}^{\mathrm{irr}})\) is
compactly generated by the set
\[
\mathcal{G}^{\mathrm{irr}}
\;=\;\bigl\{\,\iota_{x,*}\mathcal{O}_{\mathcal{G}_{x}}\otimes\det^{\otimes n}\;:\;x\in\mathcal{X}^{\mathrm{irr}}(\overline{\mathbb{F}}_{p}),\;n\in\mathbb{Z}\,\bigr\},
\]
where \(\iota_{x}\colon\mathcal{G}_{x}\hookrightarrow\mathcal{X}^{\mathrm{irr}}\)
is the residual gerbe inclusion and \(\det\) is the determinant Serre twist.
\end{proposition}

\begin{proof}
Under the QCA and closed-gerbe generation hypotheses, the
Arinkin--Gaitsgory--Gaitsgory--Rozenblyum compact-generation theorem
applies to the stratum. The determinant twists record the connected
components of the fixed-determinant gerbes. No assertion is made that
Emerton--Gee alone supplies these hypotheses at \(p=2,3\).
\end{proof}

\subsection{The Reducible Stratum}

\begin{proposition}[compact generator on the reducible stratum]
\label{v4-arith:w5-g:f1a3-C:prop:compact-gen-red}
\ClaimStatusProvedHereConditional \emph{(conditional on the two-term
local model and the gerbe-plus-extension generator theorem).} On the
stratum \(\mathcal{X}^{\mathrm{red}}\)
the category \(\mathrm{IndCoh}_{\mathrm{Nilp}}(\mathcal{X}^{\mathrm{red}})\)
is compactly generated by the set
\[
\mathcal{G}^{\mathrm{red}}
\;=\;\bigl\{\,\iota_{x,*}\mathcal{O}_{\mathcal{G}_{x}}\otimes\det^{\otimes n}\,\bigr\}_{x,n}
\;\cup\;\bigl\{\,\iota_{x,*}\mathcal{E}_{x}\otimes\det^{\otimes n}\,\bigr\}_{x,n},
\]
where \(\mathcal{E}_{x}\) is the non-split extension sheaf witnessing the
Ext\({}^{1}\)-class
\(\mathrm{Ext}^{1}(\mathbf{1}_{\mathrm{triv}},\mathbf{1}_{\mathrm{sgn}})\cong k\)
at the closed point (the essential non-semisimple data of
\Cref{v4-arith:bzn-rcompletion:prop:R4-Ext1-p2}).
\end{proposition}

\begin{proof}
The proof is formal from the stated local model. A two-term cotangent
model gives the indicated extension sheaf, and the bar resolution gives
generation by the gerbe and this extension. The equality of dimensions
does not identify the deformation class with the Bushnell--Henniart
Ext-class; that identification is part of the hypothesis.
\end{proof}

\subsection{The Wild Strata at \(p=2\) and \(p=3\)}

\begin{proposition}[compact generator on the wild stratum at \(p=2\)]
\label{v4-arith:w5-g:f1a3-C:prop:compact-gen-wild-p2}
\ClaimStatusProvedHereConditional \emph{(conditional on the \(p=2\)
wild local model and generator theorem).} At \(p=2\), the category
\(\mathrm{IndCoh}_{\mathrm{Nilp}}(\mathcal{X}^{\mathrm{wild},2})\) is
compactly generated by the set
\[
\mathcal{G}^{\mathrm{wild},2}
\;=\;\bigl\{\,\iota_{x,*}\mathcal{O}_{\mathcal{G}_{x}}\otimes\det^{\otimes n}\,\bigr\}_{x,n}
\;\cup\;\bigl\{\,\iota_{x,*}\mathcal{F}^{\mathrm{irreg}}_{x}\otimes\det^{\otimes n}\,\bigr\}_{x,n},
\]
where \(\mathcal{F}^{\mathrm{irreg}}_{x}\) is the structure sheaf of the
irregular-supersingular cycle at \(x\), corresponding under
Koszul duality to the sign-character-defect of
\Cref{v4-arith:bzn-rcompletion:prop:R4-Ext1-p2}.
\end{proposition}

\begin{proof}
Under the \(p=2\) local-model hypothesis the stratum has the same
two-term form as the reducible stratum, with one additional irregular
cycle. The bar-resolution argument then gives the displayed generator.
Bushnell--Henniart indexes types; it does not by itself prove the
ind-coherent generator statement on the Emerton--Gee stack.
\end{proof}

\begin{proposition}[compact generator on the wild strata at \(p=3\)]
\label{v4-arith:w5-g:f1a3-C:prop:compact-gen-wild-p3}
\ClaimStatusProvedHereConditional \emph{(conditional on the \(p=3\)
\(A_{2}\)-local model and generator theorem).} At \(p=3\), for each
\(i\in\{0,1,2\}\) the
category
\(\mathrm{IndCoh}_{\mathrm{Nilp}}(\mathcal{X}^{\mathrm{wild},3}_{L_{i}})\)
is compactly generated by the set
\[
\mathcal{G}^{\mathrm{wild},3}_{L_{i}}
\;=\;\bigl\{\,\iota_{x,*}\mathcal{O}_{\mathcal{G}_{x}}\otimes\det^{\otimes n}\,\bigr\}_{x,n}
\;\cup\;\bigl\{\,\iota_{x,*}\mathcal{F}^{A_{2},i}_{x}\otimes\det^{\otimes n}\,\bigr\}_{x,n},
\]
where \(\mathcal{F}^{A_{2},i}_{x}\) is the structure sheaf of the
\(A_{2}\)-quiver-edge sheaf at \(x\) (two-dimensional support, recording
the Ext\({}^{1}\)-extensions
\(\alpha_{01}\in\mathrm{Ext}^{1}(L_{1},L_{0})\) and
\(\alpha_{12}\in\mathrm{Ext}^{1}(L_{2},L_{1})\) of
\Cref{v4-arith:bzn-rcompletion:prop:R4-Ext1-p3}). The quiver-structure
relation \(\alpha_{01}\alpha_{12}=0\) ensures no further generator is
needed.
\end{proposition}

\begin{proof}
Under the \(A_{2}\)-local-model hypothesis the quiver bar resolution
has the displayed two edge generators. The relation
\(\alpha_{01}\alpha_{12}=0\) then prevents a third generator. Dimension
equality alone is not a proof of this local model.
\end{proof}

\section{Compact Generation Theorem}
\label{v4-arith:w5-g:f1a3-C:compact-gen-theorem}

\begin{theorem}[conditional compact generation of
\(\mathrm{IndCoh}_{\mathrm{Nilp}}(\mathcal{X}^{\mathrm{EG},p,\chi}_{GL_{2}})\)
at wild primes]
\label{v4-arith:w5-g:f1a3-C:thm:main}
\ClaimStatusProvedHereConditional. Fix \(p\in\{2,3\}\) and a central character
\(\chi\colon\mathbb{Z}_{p}^{\times}\to k^{\times}\). The \(\infty\)-category
\(\mathrm{IndCoh}_{\mathrm{Nilp}}(\mathcal{X}^{\mathrm{EG},p,\chi}_{GL_{2}})\)
is compactly generated if the finite Nilp stratification of
\Cref{v4-arith:w5-g:f1a3-C:prop:nilp-stratification} and the
per-stratum generator statements of
\Cref{v4-arith:w5-g:f1a3-C:prop:compact-gen-irr},
\Cref{v4-arith:w5-g:f1a3-C:prop:compact-gen-red},
\Cref{v4-arith:w5-g:f1a3-C:prop:compact-gen-wild-p2}, and
\Cref{v4-arith:w5-g:f1a3-C:prop:compact-gen-wild-p3}
hold. Under these inputs, explicit compact generators are given by the
union across strata:
\begin{equation}
\label{v4-arith:w5-g:f1a3-C:eq:full-compact-generator}
\mathcal{G}^{\mathrm{full},p}
\;=\;\mathcal{G}^{\mathrm{irr}}\;\cup\;\mathcal{G}^{\mathrm{red}}\;\cup\;\mathcal{G}^{\mathrm{wild},p},
\end{equation}
where \(\mathcal{G}^{\mathrm{wild},p}\) denotes the full wild generator
set of \Cref{v4-arith:w5-g:f1a3-C:prop:compact-gen-wild-p2} (at \(p=2\))
or the union \(\bigcup_{i}\mathcal{G}^{\mathrm{wild},3}_{L_{i}}\) of
\Cref{v4-arith:w5-g:f1a3-C:prop:compact-gen-wild-p3} (at \(p=3\)).
\end{theorem}

\begin{proof}
Assume the stated stratification and per-stratum generator inputs.
Recollement for a finite locally closed decomposition reduces compact
generation of the total category to compact generation on the strata.
The compact generators on the strata push forward along the closed
pieces and extend along the open pieces. Induction over the finite
stratification gives the displayed generator set.
\end{proof}

\begin{corollary}[Wild-prime compact generation and the Bushnell--Henniart indexing map]
\label{v4-arith:w5-g:f1a3-C:cor:F1a3C-discharged}
\ClaimStatusProvedHereConditional \emph{(on the class-level
identification of Bushnell--Henniart closed points with residual
gerbes).} The compact-generation part of the conjecture
\Cref{v4-arith:f1a3-BH:conj:wild-compact-generation} (F1.a.3.C) holds
at \(p\in\{2,3\}\) for the classical Emerton--Gee stack if the
finite Nilp stratification and per-stratum generator inputs hold: the
category
\(\mathrm{IndCoh}_{\mathrm{Nilp}}(\mathcal{X}^{\mathrm{EG},p,\chi}_{GL_{2}})\)
is compactly generated by explicit residual-gerbe plus
Serre-twist generators. The Bushnell--Henniart closed points
\(x_{\tau}\) of Chapter~\ref{v4-arith:f1a3-BH} give the indexing map to
codimension-zero strata; identifying these points with residual gerbes
is a class-level statement, not a consequence of the equality of
dimensions alone.
\end{corollary}

\begin{proof}
Conditional compact generation is
\Cref{v4-arith:w5-g:f1a3-C:thm:main}. The identification of BH closed
points with residual gerbes of each stratum requires a matching of indexing
sets: BH types of \Cref{v4-arith:f1a3-BH:prop:BH-types-p2} and
\Cref{v4-arith:f1a3-BH:prop:BH-types-p3} correspond bijectively to
residual Galois representations on the strata of
\Cref{v4-arith:w5-g:f1a3-C:prop:nilp-stratification} (by BH local
Langlands \S 18 for \(GL_{2}(\mathbb{Q}_{p})\) + Emerton--Gee \S 5.4
for the spectral-side decomposition). This comparison is retained as
the named conditional input; dimension agreement is not used in the
reverse direction.
\end{proof}

\begin{remark}[independence from the wild-EG enhancement conjecture]
\label{v4-arith:w5-g:f1a3-C:rem:independence}
\Cref{v4-arith:w5-g:f1a3-C:thm:main} and its corollary
\Cref{v4-arith:w5-g:f1a3-C:cor:F1a3C-discharged} do \emph{not} invoke
\Cref{v4-arith:deg-f1a:conj:EG-wild-enhancement} (F1.a.2, the wild-EG
derived enhancement). They are conditional statements about the
classical Emerton--Gee stack and its \(\mathrm{IndCoh}_{\mathrm{Nilp}}\).
The derived enhancement would be a further input, not a consequence of
the classical compact-generation implication.
\end{remark}

\section{Residual Obstruction Statement}
\label{v4-arith:w5-g:f1a3-C:residual-obstruction}

\subsection{Domain of the Classical Result}

\begin{proposition}[domain of F1.a.3.C]
\label{v4-arith:w5-g:f1a3-C:prop:domain-summary}
\ClaimStatusProvedHere. The F1.a.3.C decomposition consists of:
\begin{enumerate}
\item \(\mathcal{X}^{\mathrm{EG},p,\chi}_{GL_{2}}\) is locally
Noetherian and quasi-excellent at every prime \(p\)
(\Cref{v4-arith:w5-g:f1a3-C:prop:EG-noetherian},
\Cref{v4-arith:w5-g:f1a3-C:prop:quasi-excellent}); this is the
unconditional geometric base.
\item The nilpotent-cone locus decomposes into a finite set of
locally closed Noetherian substacks, with the irregular supersingular
component at \(p=2\) and the \(A_{2}\)-quiver-stratified principal block
at \(p=3\)
(\Cref{v4-arith:w5-g:f1a3-C:prop:nilp-stratification}); this is a
conditional wild-prime local-model input.
\item Each stratum admits explicit compact generators: residual-gerbe
structure sheaves tensored with Serre twists of the determinant character,
together with the single or double extension sheaves encoding the wild
\(\mathrm{Ext}^{1}\)-data
(\Cref{v4-arith:w5-g:f1a3-C:prop:compact-gen-irr},
\Cref{v4-arith:w5-g:f1a3-C:prop:compact-gen-red},
\Cref{v4-arith:w5-g:f1a3-C:prop:compact-gen-wild-p2},
\Cref{v4-arith:w5-g:f1a3-C:prop:compact-gen-wild-p3}); this is
conditional on the corresponding gerbe-generation theorem.
\item The full category
\(\mathrm{IndCoh}_{\mathrm{Nilp}}(\mathcal{X}^{\mathrm{EG},p,\chi}_{GL_{2}})\)
is compactly generated at \(p\in\{2,3\}\)
(\Cref{v4-arith:w5-g:f1a3-C:thm:main}); this is conditional on the
previous two inputs.
\item The conjecture of Chapter~\ref{v4-arith:f1a3-BH} holds in its
classical-stack form
(\Cref{v4-arith:w5-g:f1a3-C:cor:F1a3C-discharged}); this is conditional,
not an unconditional closure.
\end{enumerate}
\end{proposition}

\subsection{Domain Boundary: Derived Enhancement}

\begin{remark}[domain boundary: derived wild-EG enhancement]
\label{v4-arith:w5-g:f1a3-C:residual:derived-extension}
\Cref{v4-arith:deg-f1a:conj:EG-wild-enhancement} remains open.
Compact generation on the classical stack does not imply compact
generation on the conjectural derived formal stack
\(\mathcal{X}^{\mathrm{EG,wild},p,\chi}_{GL_{2}}\); upgrading requires F1.a.2.
Chapter~\ref{v4-arith:f1a3-BH} demands the classical statement; this
chapter reduces it to the finite-stratification and gerbe-generation
inputs. The derived enhancement is a separate conjecture of
Chapter~\ref{v4-arith:deg-f1a}.
\end{remark}

\begin{remark}[finite Nilp stratification is an input]
\label{v4-arith:w5-g:f1a3-C:rem:no-boundary-obstruction}
The stratification of
\Cref{v4-arith:w5-g:f1a3-C:prop:nilp-stratification} is exhaustive only
under the stated wild-prime local-model hypothesis:
the three strata at \(p=2\) or five strata at \(p=3\) cover the full
nilpotent cone. Emerton--Gee gives the ambient stack and deformation
theory; it does not by itself exclude additional Nilp components at
higher codimension.
\end{remark}

\begin{remark}[source disjointness]
\label{v4-arith:w5-g:f1a3-C:rem:source-disjointness}
The proof draws on five independent source families.
(1) Emerton--Gee arXiv:1908.07185 supplies the Noetherian formal
algebraic stack. Bhatt--Scholze arXiv:1905.08229 supplies the
excellent formal-neighbourhood input. Dotto--Emerton--Gee
arXiv:2603.26887 is a \(p\geq 5\) theorem and supplies no wild-prime
closure.
(2) Gaitsgory--Rozenblyum DAG Vol.~II provides compact-generator transfer
and recollement.
(3) Arinkin--Gaitsgory arXiv:1201.6343 provides the nilpotent-singular-support
framework.
(4) Bushnell--Henniart 2006 provides the wild-type classification.
(5) Chapter~\ref{v4-arith:bzn-rcompletion} provides the \(\mathrm{Ext}^{1}\)-computation.
These five families are mutually compatible: Emerton--Gee uses the
\((\varphi,\Gamma)\)-module conventions of Fontaine--Wintenberger; the
Nilp convention follows Arinkin--Gaitsgory (singular support as a closed
conic subset of the cotangent stack); the Serre twist convention of
Gaitsgory--Rozenblyum DAG Vol.~II Ch.~5 \S 5.4 matches the
Emerton--Gee determinant stratification. All five families are disjoint
from the bar--cobar corpus of Vols.~I--III and from the
Fargues--Scholze prismatic classification. Recollement across strata
is functorial; closed immersions preserve compact generators (DAG
Vol.~II Ch.~4 Prop.~4.4). The hypotheses required beyond the ambient
Noetherian stack are the finite wild-prime Nilp stratification and the
per-stratum generator theorem. These are not supplied by Emerton--Gee
alone. No numerical computation is involved.
\end{remark}

\section{Bibliography Stanza}
\label{v4-arith:w5-g:f1a3-C:biblio}

The primary sources invoked, with disjointness rationale:

\begin{description}
\item[Emerton--Gee arXiv:1908.07185.] The Noetherian formal algebraic
structure of \(\mathcal{X}^{\mathrm{EG},p,\chi}_{GL_{2}}\) at every
prime \(p\geq 2\). Used in \S\ref{v4-arith:w5-g:f1a3-C:locally-noetherian}
(Noetherian structure). It does not prove the finite Nilp
stratification used later.
\item[Emerton--Gee arXiv:1908.07185.] The moduli-of-\((\varphi,\Gamma)\)-modules
construction and tangent-complex analysis. Used as ambient
deformation theory, not as an explicit wild-prime compact-generator
theorem.
\item[Dotto--Emerton--Gee arXiv:2603.26887.] The tame-DEG machinery
\(p\geq 5\). It is evidence for the shape of the functor at tame
primes. It is not an input proving compact generation at \(p=2,3\).
\item[Bhatt--Scholze arXiv:1905.08229.] The prism theory supplying
quasi-excellence of formal neighbourhoods on Noetherian excellent
rings; Proposition 3.8 quoted in
\Cref{v4-arith:w5-g:f1a3-C:prop:quasi-excellent}.
\item[Gaitsgory--Rozenblyum AMS DAG Vol.~II.] The locally-Noetherian
compact generation of \(\mathrm{IndCoh}\) (Ch.~4 Prop.~4.4); recollement
across strata (Ch.~II \S 2); Serre-twist filling (Ch.~5 \S 5.4);
bar-resolution compact generation on tor-amplitude-\([-1,0]\) stacks
(Ch.~8 \S 8.3). Used throughout \S\ref{v4-arith:w5-g:f1a3-C:compact-generators}
and \S\ref{v4-arith:w5-g:f1a3-C:compact-gen-theorem}.
\item[Arinkin--Gaitsgory arXiv:1201.6343.] Nilpotent singular support;
Theorem 12.5 for compact generation of
\(\mathrm{IndCoh}_{\mathrm{Nilp}}\) by closed-point structure sheaves on
smooth locally Noetherian stacks. Used in
\Cref{v4-arith:w5-g:f1a3-C:prop:compact-gen-irr} Step~4.
\item[Bushnell--Henniart 2006.] Local Langlands and wild-type
classification. It supplies representation-theoretic indexing data.
It does not identify the Nilp components of the Emerton--Gee stack.
\item[Beilinson--Bernstein--Deligne 1982.] Recollement formalism;
extended to DAG via Gaitsgory--Rozenblyum. Used in
\Cref{v4-arith:w5-g:f1a3-C:thm:main} Step~1.
\item[Beilinson--Ginzburg 1996.] Koszul duality between cotangent fibres
and Yoneda Ext-algebras. It explains the expected shape of the local
model. The actual identification of the wild-prime cotangent directions
with the BH-block Ext\({}^{1}\)-classes remains part of the conditional
input.
\item[Chapter~\ref{v4-arith:bzn-rcompletion}.] The Ext\({}^{1}\)-data at
\(p=2\) (\Cref{v4-arith:bzn-rcompletion:prop:R4-Ext1-p2}) and \(p=3\)
(\Cref{v4-arith:bzn-rcompletion:prop:R4-Ext1-p3}) supplying the
expected dimensions of the wild-prime Ext directions. Equal dimensions
do not prove equality of local models.
\end{description}

All sources are disjoint from the proof sources of the Vol.~IV
realization functor \(\mathrm{Real}\) (Pridham--Toen--Vezzosi; Lurie HA
5.3; Costello--Gwilliam; Francis--Gaitsgory): none of these appear in
the proof. The formal recollement implication is self-contained within
the Emerton--Gee + Gaitsgory--Rozenblyum + Arinkin--Gaitsgory corpus.
The wild-prime local models are separate inputs.

\begin{remark}[consequence for the F1.a obstruction decomposition]
\label{v4-arith:w5-g:f1a3-C:rem:F1a-reduction}
\Cref{v4-arith:w5-g:f1a3-C:cor:F1a3C-discharged} is conditional. The
F1.a obstruction decomposition of Chapter~\ref{v4-arith:f1a3-BH} therefore
retains three open items: the Bushnell--Henniart type-block comparison
conjecture \Cref{v4-arith:f1a3-BH:conj:BH-packaging}, the Paskunas
wild \(\mathrm{Ext}^{1}\)-compatibility conjecture
\Cref{v4-arith:f1a3-BH:conj:Paskunas-wild-ext}, and the wild-prime
compact-generation local model of this chapter.
\end{remark}
